\section{Kiểm thử Cơ sở dữ liệu}

\subsection{Phạm vi kiểm thử}
Vì hệ thống chưa có Backend thật (mục~\ref{sec:scope-limitations}), kiểm thử trong phạm vi đồ án diễn ra hoàn toàn ở \textbf{mức cơ sở dữ liệu} — chạy trực tiếp script T-SQL trên SQL Server Management Studio (SSMS) để kiểm chứng các ràng buộc và công thức nghiệp vụ trọng yếu:
\begin{itemize}
    \item Ràng buộc append-only trên \texttt{inventory\_transactions} (thử UPDATE/DELETE, kỳ vọng bị chặn bởi INSTEAD OF TRIGGER).
    \item Công thức giá vốn bình quân gia quyền (so sánh giá trị \texttt{avg\_cost} tính toán với giá trị kỳ vọng theo kịch bản nhập hàng nhiều đợt ở mục~4.2.7).
    \item Ràng buộc \texttt{available = on\_hand - reserved} không âm (đã được ép buộc ngay ở tầng schema bằng \texttt{CHECK} constraint, xem mục~4.2.3).
    \item Tính duy nhất của \texttt{sku\_code} trong phạm vi Tenant (BR-03) và tính duy nhất toàn hệ thống của \texttt{barcodes.code}.
    \item Cô lập dữ liệu đa tenant: truy vấn chéo TenantId phải trả về rỗng.
    \item Đối chiếu tổng Ledger với StockBalance (kiểm chứng tính đúng đắn của "read model" đã phân tích ở mục~4.2.2).
\end{itemize}

\textbf{Không nằm trong phạm vi kiểm thử của học kỳ này:} Kiểm thử tải thực tế qua API (hiện chỉ mô phỏng bằng nhiều phiên SSMS song song như mục~4.3.7), kiểm thử giao diện (chưa có Frontend), kiểm thử tích hợp Webhook thật.

\subsection{Kiểm chứng ràng buộc Append-only}
Chạy trực tiếp 2 câu lệnh dưới đây trên một dòng \texttt{inventory\_transactions} đã tồn tại từ seed data, cả hai đều \textbf{phải thất bại} (hoặc chính xác hơn: không có tác dụng gì, vì bị \texttt{INSTEAD OF TRIGGER} thay thế):

\begin{lstlisting}
-- Test 1: thu UPDATE truc tiep
UPDATE inventory_transactions
SET quantity = 999
WHERE id = @existing_transaction_id;
-- Ket qua ky vong:
-- Msg 50000, Level 16, State 1
-- inventory_transactions is append-only. UPDATE/DELETE is not allowed.
-- Use a compensating entry instead.
-- (0 rows affected) -- du lieu KHONG bi thay doi

-- Test 2: thu DELETE truc tiep
DELETE FROM inventory_transactions
WHERE id = @existing_transaction_id;
-- Ket qua ky vong: cung thong bao loi tren, (0 rows affected)
\end{lstlisting}

Cả hai câu lệnh đều bị trigger \texttt{trg\_block\_ledger\_modification} (mục~4.2.4) chặn lại — xác nhận NFR-SEC-03 được thực thi đúng ở tầng cơ sở dữ liệu, không phụ thuộc vào việc Backend có kiểm tra hay không.

\subsection{Kiểm chứng cô lập dữ liệu đa Tenant}
Truy vấn đối chiếu \texttt{tenant\_id} giữa bảng cha-con — kỳ vọng \textbf{0 dòng} trả về (không có bản ghi nào bị lệch Tenant do lỗi ứng dụng):
\begin{lstlisting}
SELECT o.order_id, o.tenant_id AS order_tenant, b.tenant_id AS branch_tenant
FROM orders o JOIN branches b ON o.branch_id = b.branch_id
WHERE o.tenant_id <> b.tenant_id;
-- Ket qua ky vong: 0 dong

SELECT sb.branch_id, sb.sku_id, sb.tenant_id AS balance_tenant, sv.tenant_id AS sku_tenant
FROM stock_balances sb JOIN sku_variants sv ON sb.sku_id = sv.sku_id
WHERE sb.tenant_id <> sv.tenant_id;
-- Ket qua ky vong: 0 dong
\end{lstlisting}
Lặp lại truy vấn tương tự cho toàn bộ 9 bảng con có \texttt{tenant\_id} denormalized (mục~4.2.2, Ngoại lệ 1) để đảm bảo tính nhất quán được duy trì trong suốt vòng đời dữ liệu.


\subsection{Kiểm chứng tính toàn vẹn Ledger $\leftrightarrow$ StockBalance}
Script dưới đây đối chiếu số dư trong \texttt{stock\_balances} với tổng cộng dồn từ toàn bộ lịch sử \texttt{inventory\_transactions} — nếu thiết kế và dữ liệu đúng, truy vấn phải trả về \textbf{0 dòng} (không có sai lệch):

\begin{lstlisting}
SELECT
    sb.sku_id,
    sb.branch_id,
    sb.on_hand AS balance_on_hand,
    SUM(CASE WHEN it.transaction_type = 'IN' THEN it.quantity
             ELSE -it.quantity END)   AS ledger_computed_on_hand
FROM stock_balances sb
JOIN inventory_transactions it
    ON it.sku_id = sb.sku_id AND it.branch_id = sb.branch_id
GROUP BY sb.sku_id, sb.branch_id, sb.on_hand
HAVING sb.on_hand <>
    SUM(CASE WHEN it.transaction_type = 'IN' THEN it.quantity
             ELSE -it.quantity END);
-- Ket qua ky vong: 0 dong (khong co SKU nao bi lech)
\end{lstlisting}

\subsection{Kiểm chứng tổng tiền đơn hàng (Orders.TotalAmount)}
Đối chiếu cột cache \texttt{orders.total\_amount} (mục~4.2.2, Ngoại lệ 2) với tổng thực tế từ \texttt{order\_items.subtotal} — kỳ vọng \textbf{0 dòng} lệch:
\begin{lstlisting}
SELECT o.order_id, o.status, o.total_amount,
       SUM(oi.subtotal) AS item_total,
       o.total_amount - SUM(oi.subtotal) AS difference
FROM orders o
JOIN order_items oi ON oi.order_id = o.order_id
GROUP BY o.order_id, o.status, o.total_amount
HAVING o.total_amount <> SUM(oi.subtotal);
-- Ket qua ky vong: 0 dong
\end{lstlisting}

\subsection{Kiểm chứng công thức giá vốn bình quân}
Dựa trên kịch bản seed data 2 đợt nhập tại mục~4.2.7 (100 đơn vị giá 50.000đ, sau đó thêm 50 đơn vị giá 56.000đ), giá trị \texttt{avg\_cost} kỳ vọng được tính tay như sau:
\[
\text{New Cost} = \frac{(100 \times 50.000) + (50 \times 56.000)}{100 + 50} = \frac{5.000.000 + 2.800.000}{150} = 52.000\ \text{đ}
\]
Script kiểm tra:
\begin{lstlisting}
SELECT sku_id, branch_id, avg_cost
FROM stock_balances
WHERE sku_id = @sku_id AND branch_id = @branch_id;
-- Ket qua ky vong: avg_cost = 52000.0000
\end{lstlisting}

\subsection{Mô phỏng Concurrency (kịch bản Flash Sale)}
\label{sec:concurrency-test}
Để kiểm chứng cơ chế ở mục~\ref{sec:concurrency-control} tại mức thiết kế (chưa có Backend để test tải qua API thật), nhóm mở 2 cửa sổ Query trong SSMS (mô phỏng 2 phiên kết nối song song), mỗi phiên mô phỏng một request Reserve trên cùng một SKU chỉ còn 1 đơn vị tồn:

\begin{lstlisting}
-- Cua so A (chay truoc, giu khoa):
BEGIN TRANSACTION;
SELECT on_hand, reserved FROM stock_balances WITH (UPDLOCK, ROWLOCK)
WHERE sku_id = 'sku-flash-sale';
-- (chua COMMIT, giu khoa)

-- Cua so B (chay dong thoi, cung SKU):
BEGIN TRANSACTION;
SELECT on_hand, reserved FROM stock_balances WITH (UPDLOCK, ROWLOCK)
WHERE sku_id = 'sku-flash-sale';
-- Cua so B bi treo (block), phai cho toi khi Cua so A COMMIT/ROLLBACK
\end{lstlisting}

\textbf{Quan sát kỳ vọng:} Cửa sổ B phải ở trạng thái chờ (hiển thị "Executing query..." không trả kết quả) cho đến khi Cửa sổ A thực hiện \texttt{COMMIT TRANSACTION} hoặc \texttt{ROLLBACK TRANSACTION} — xác nhận SQL Server tuần tự hóa đúng 2 giao dịch tranh chấp trên cùng 1 dòng, không có hiện tượng 2 giao dịch cùng đọc được \texttt{reserved} cũ và cùng ghi đè (lost update). Có thể quan sát trực quan hơn qua view hệ thống \texttt{sys.dm\_tran\_locks} để thấy Cửa sổ B đang ở trạng thái \texttt{WAIT} chờ khóa \texttt{UPDLOCK} mà Cửa sổ A đang giữ.

\subsection{Kết quả kiểm thử}

\begin{table}[htbp]
    \centering
    \small
    \caption{Tổng hợp kết quả kiểm thử mức Database}
    \begin{tabularx}{\textwidth}{|Y|p{2.2cm}|p{1.6cm}|p{1.6cm}|Y|}
        \hline
        \textbf{Hạng mục kiểm thử} & \textbf{Số test case} & \textbf{Pass} & \textbf{Fail} & \textbf{Ghi chú} \\
        \hline
        Schema Validation (DDL chạy sạch trên container mới) & -- & -- & -- & Điền số liệu thực tế sau khi chạy \\
        \hline
        Append-only Constraint (mục 4.3.2) & -- & -- & -- & \\
        \hline
        Cô lập dữ liệu đa Tenant (mục 4.3.3) & -- & -- & -- & \\
        \hline
        Toàn vẹn Ledger $\leftrightarrow$ StockBalance (mục 4.3.4) & -- & -- & -- & \\
        \hline
        Tổng tiền đơn hàng (mục 4.3.5) & -- & -- & -- & \\
        \hline
        Công thức giá vốn bình quân (mục 4.3.6) & -- & -- & -- & So sánh với 52.000đ tính tay \\
        \hline
        Ràng buộc available không âm (CHECK constraint) & -- & -- & -- & \\
        \hline
        Concurrency simulation (mục 4.3.7) & -- & -- & -- & Quan sát hành vi khóa qua sys.dm\_tran\_locks \\
        \hline
    \end{tabularx}
\end{table}

\textit{Ghi chú: Bảng trên để trống cột số liệu — nhóm điền kết quả thật sau khi chạy xong toàn bộ script kiểm thử trên môi trường Docker SQL Server, đính kèm chi tiết đầy đủ tại file kết quả kiểm thử theo quy ước bàn giao ở Chương~2.}

\subsection{Lỗi phát hiện và hướng xử lý}
Các lỗi phát hiện trong quá trình kiểm thử schema/dữ liệu được ghi nhận trên GitHub Issues theo mức độ ưu tiên (Critical/High/Medium/Low), tuân theo quy trình Quản lý chất lượng đã mô tả ở Chương~2. Lỗi liên quan đến ràng buộc append-only hoặc công thức giá vốn được xếp mức Critical vì ảnh hưởng trực tiếp đến tính đúng đắn của toàn bộ hệ thống, ưu tiên xử lý trong vòng 24 giờ.
